\Chapter{El equilibrio de Wardrop}

\section{Definición y propiedades del equilibrio}

\subsection{Los principios de Wardrop}

\datebox{2026}

En su artículo seminal de 1952, John Glen Wardrop formuló dos principios que capturan los posibles objetivos en la asignación de tráfico:

\begin{definition}[Principios de Wardrop (1952)]\label{def:wardrop_principles}
\begin{enumerate}
    \item \textbf{Principio I (Equilibrio de usuario, UE):} \emph{``Las condiciones de flujo en la red son tales que ningún viajero individual puede reducir su tiempo de viaje cambiando unilateralmente de ruta.''}
    \item \textbf{Principio II (Óptimo social, SO):} \emph{``El tiempo de viaje promedio por vehículo es mínimo.''}
\end{enumerate}
\end{definition}

\begin{remark}
El Principio I de Wardrop es una condición de \emph{estabilidad}: ningún conductor tiene incentivo para desviarse. El Principio II es una condición de \emph{optimalidad colectiva}. En general, las asignaciones que satisfacen uno no satisfacen el otro, y la brecha entre ambas es el objeto de estudio central de estas notas.
\end{remark}

\subsection{Definición formal del equilibrio de Wardrop}

\begin{definition}[Equilibrio de Wardrop]\label{def:wardrop_equilibrium}
Una asignación $f^* \in \FeasSet$ es un \emph{equilibrio de Wardrop} (o equilibrio de usuario) si para todo par $w \in W$ existe un escalar $\mu_w^* \in \R$ tal que:
\begin{align}
f_p^* > 0 &\implies C_p(x^*) = \mu_w^*, \label{eq:wardrop_used}\\
f_p^* = 0 &\implies C_p(x^*) \geq \mu_w^*, \label{eq:wardrop_unused}
\end{align}
donde $x^* = \Delta f^*$. El escalar $\mu_w^*$ se denomina el \emph{costo de equilibrio} del par $w$.
\end{definition}

\begin{remark}
Las condiciones~\upeqref{eq:wardrop_used}--\upeqref{eq:wardrop_unused} expresan que: (i) todas las rutas que llevan flujo positivo tienen el mismo costo, y (ii) ese costo es mínimo entre todas las rutas disponibles. En términos prácticos: en el equilibrio, todos los conductores de un mismo par O-D experimentan el mismo tiempo de viaje, y no existe otra ruta que les tomara menos tiempo.
\end{remark}

\begin{definition}[Sistema de complementariedad de Wardrop]
Las condiciones~\upeqref{eq:wardrop_used}--\upeqref{eq:wardrop_unused} se escriben equivalentemente como el siguiente sistema de complementariedad no lineal: para todo $w \in W$ y todo $p \in \calP_w$,
\begin{equation}\label{eq:complementarity}
f_p^* \geq 0, \quad C_p(x^*) - \mu_w^* \geq 0, \quad f_p^* \cdot \left( C_p(x^*) - \mu_w^* \right) = 0.
\end{equation}
\end{definition}

\subsection{Conexión con el equilibrio de Nash}

\begin{definition}[Juego de congestión en forma normal]
Un \emph{juego de congestión} con $N$ jugadores es una tupla $\Gamma = (I, (S_i)_{i \in I}, (C_i)_{i \in I})$ donde:
\begin{itemize}
    \item $I = \{1, \ldots, N\}$ es el conjunto de jugadores (conductores);
    \item $S_i = \calP_{w(i)}$ es el espacio de estrategias del jugador $i$ (las rutas disponibles para su par O-D $w(i)$);
    \item $C_i(p_i, p_{-i}) = C_{p_i}(x(p))$ es el costo del jugador $i$ cuando se usa el perfil de estrategias $p = (p_i, p_{-i})$, donde $x_a(p) = |\{j : a \in p_j\}|/N$ es el flujo normalizado.
\end{itemize}
\end{definition}

\begin{definition}[Equilibrio de Nash]
Un perfil de estrategias $p^* = (p_1^*, \ldots, p_N^*)$ es un \emph{equilibrio de Nash} si para todo $i \in I$ y toda estrategia alternativa $q \in S_i$:
\[
C_i(p_i^*, p_{-i}^*) \leq C_i(q, p_{-i}^*).
\]
\end{definition}

\begin{theorem}[Wardrop como límite de Nash]\label{thm:wardrop_nash}
Sea $\{(\Gamma^N, f^{*,N})\}_{N \geq 1}$ una secuencia de juegos de congestión donde $N \to \infty$, cada jugador controla una fracción $1/N$ de la demanda total, y $f^{*,N}$ es un equilibrio de Nash de $\Gamma^N$. Entonces, toda acumulación de $\{f^{*,N}\}$ es un equilibrio de Wardrop.
\end{theorem}

\begin{proof}[Esquema de la demostración]
En el equilibrio de Nash con $N$ jugadores, el conductor $i$ que usa la ruta $p^*_i$ satisface:
\[
C_{p^*_i}\left(x^{*,N} + \frac{1}{N} e_{p^*_i}\right) \leq C_q\left(x^{*,N} - \frac{1}{N} e_{p^*_i} + \frac{1}{N} e_q\right) \quad \forall q \in S_i,
\]
donde $e_p$ denota el vector canónico de la ruta $p$. Cuando $N \to \infty$, la perturbación $1/N$ tiende a cero. Por continuidad de $t_a$ y por compacidad del conjunto factible, cualquier punto de acumulación $f^*$ satisface:
\[
C_{p}(x^*) \leq C_q(x^*) \quad \forall p \text{ con } f_p^* > 0, \forall q \in \calP_w,
\]
que es exactamente la condición de Wardrop~\upeqref{eq:wardrop_used}--\upeqref{eq:wardrop_unused}. Para la demostración completa véase \cite{Roughgarden2005}, Capítulo 2.
\end{proof}

\begin{remark}
El Teorema~\upref{thm:wardrop_nash} da la justificación formal del modelo de Wardrop: cuando hay ``muchos'' conductores, cada uno de los cuales es infinitesimalmente pequeño en relación al flujo total, su elección racional individual conduce al equilibrio de Wardrop. Este es el supuesto estándar en el análisis de redes viales urbanas, donde el número de conductores es del orden de los miles.
\end{remark}

\section{Existencia y unicidad del equilibrio}

\subsection{Existencia}

\begin{theorem}[Existencia del equilibrio de Wardrop]\label{thm:existence}
Si $t_a : \R_+ \to \R_+$ es continua para todo $a \in A$, y si $\FeasSet \neq \emptyset$, entonces existe al menos un equilibrio de Wardrop.
\end{theorem}

\begin{proof}
Definimos el operador de \emph{mejor respuesta}:
\[
\Phi(f) \defeq \arg\min_{g \in \FeasSet} \sum_{a \in A} t_a(x_a(f)) \cdot y_a(g),
\]
donde $y_a(g) = \sum_p \delta_{ap} g_p$. Este operador asigna a cada asignación $f$ el flujo que minimiza el costo total dado los tiempos de viaje inducidos por $f$. Nótese que $\Phi(f)$ corresponde a enrutar toda la demanda de cada par O-D por la ruta más corta dada $f$.

El operador $\Phi$ mapea el conjunto compacto y convexo $\FeasSet$ en sí mismo. Por continuidad de $t_a$ y del costo de ruta $C_p$, $\Phi$ es una correspondencia de conjuntos convexos con grafo cerrado. Por el \emph{teorema de punto fijo de Kakutani}, existe $f^* \in \Phi(f^*)$.

Un punto fijo $f^* \in \Phi(f^*)$ satisface: para todo par $w$ y toda ruta $q \in \calP_w$,
\[
\sum_{a} t_a(x_a^*) x_a(f^*) \leq \sum_a t_a(x_a^*) y_a(g) \quad \forall g \in \FeasSet,
\]
lo que implica que $f^*$ asigna flujo solo a rutas de mínimo costo dado $x^*$. Esta es exactamente la condición de Wardrop.
\end{proof}

\subsection{Unicidad}

\begin{theorem}[Unicidad de los flujos en aristas]\label{thm:uniqueness}
Si $t_a$ es estrictamente creciente para todo $a \in A$, entonces el vector de flujos en aristas $x^* \in \R^m_+$ es único en el equilibrio de Wardrop. Los flujos por ruta $f^*$ pueden no ser únicos.
\end{theorem}

\begin{proof}
Supongamos que $f^1, f^2 \in \FeasSet$ son dos equilibrios de Wardrop con flujos en aristas $x^1 = \Delta f^1$ y $x^2 = \Delta f^2$, y supongamos que $x^1 \neq x^2$.

Sea $f^\lambda = \lambda f^1 + (1-\lambda) f^2$ para $\lambda \in (0,1)$, con flujos $x^\lambda = \lambda x^1 + (1-\lambda) x^2$. Dado que $\FeasSet$ es convexo, $f^\lambda \in \FeasSet$.

Consideremos la función objetivo de Beckmann $Z(f) = \sum_a T_a(x_a(f))$ (que definiremos formalmente en el Capítulo~\ref{cap:beckmann}). Por convexidad estricta de $T_a$ y por la hipótesis $x^1 \neq x^2$, se tiene $Z(f^\lambda) < \lambda Z(f^1) + (1-\lambda) Z(f^2)$.

Pero en el Capítulo~\ref{cap:beckmann} demostraremos (Teorema~\ref{thm:beckmann}) que tanto $f^1$ como $f^2$ minimizan $Z$ sobre $\FeasSet$, por lo que $Z(f^1) = Z(f^2) = \min_{\FeasSet} Z$. Esto contradice la desigualdad estricta anterior. Luego $x^1 = x^2$.

Para ver que los flujos por ruta no son únicos en general, considérese una red con dos aristas paralelas de igual costo: cualquier distribución del flujo entre ellas es un equilibrio.
\end{proof}

\begin{remark}
La unicidad de los flujos en aristas (no por rutas) es la propiedad relevante para la ingeniería de tráfico, pues los flujos en aristas son las cantidades observables (mediante aforadores vehiculares).
\end{remark}
